\documentclass[11pt,a4paper]{article}

\usepackage[margin=1in]{geometry}
\usepackage{amsmath,amssymb,amsthm}
\usepackage{enumitem}
\usepackage[hidelinks]{hyperref}

\newtheorem{theorem}{Theorem}[section]
\newtheorem{proposition}[theorem]{Proposition}
\newtheorem{lemma}[theorem]{Lemma}
\theoremstyle{definition}
\newtheorem{definition}[theorem]{Definition}
\theoremstyle{remark}
\newtheorem{remark}[theorem]{Remark}

\setlist[itemize]{leftmargin=*,topsep=3pt,itemsep=2pt}
\setlist[enumerate]{leftmargin=*,topsep=3pt,itemsep=2pt}

\newcommand{\Tr}{\operatorname{Tr}}
\newcommand{\id}{\mathbf{1}}
\newcommand{\cH}{\mathcal{H}}
\newcommand{\cI}{\mathcal{I}}
\newcommand{\cK}{\mathcal{K}}
\newcommand{\cD}{\mathcal{D}}
\newcommand{\cB}{\mathcal{B}}
\newcommand{\cU}{\mathcal{U}}
\newcommand{\cX}{\mathcal{X}}
\newcommand{\Pcoh}{\Pi_{\mathrm{coh}}}
\newcommand{\Var}{\operatorname{Var}}

\title{Measurement as Physical Disturbance, Outcome Completion,\\
and Record Stabilization in Modal Triplet Theory}

\author{Peter Nero}

\date{12 September 2026, Version 7}

\begin{document}
\maketitle

\begin{abstract}
Measurement is an ordinary physical interaction, but an adequate mathematical
description must distinguish three stages: apparatus coupling, completion into
an outcome-resolved transition, and stabilization of the resulting record.
This paper formulates that distinction for Modal Triplet Theory (MTT). A
localized coupling may displace a state from a quiet coherent regime, and
contractive dynamics inside a selected record basin may then stabilize a
repeatable record. Neither fact chooses a basin. We therefore introduce the
missing transition-completion object: a normalized kernel, or in quantum
language an instrument, that assigns both an outcome probability and a
post-outcome state. We prove that one global contraction cannot support
multiple stable outcomes, derive the exact normalization and conditional-state
laws of a completion kernel, and distinguish failure of unique decoding from
the possible existence of a representative section. Decoherence suppresses
interference within an outcome algebra but does not select one instrument
element. The Born rule requires a selected preparation law and a basin--trace
or instrument--trace equality. The current q79 binary one-anchor recorder
supplies such an exact stopped-output law on its declared domain; arbitrary
apparatus contexts and objective one-history selection remain open. An
Ornstein--Uhlenbeck variance is retained only as a conditional linear-response
model and is not identified with the Heisenberg uncertainty principle.
\end{abstract}

\section*{Version 7 Revision Note}
\begin{description}[style=nextline]
\item[Supersedes] Version 6 as the current manuscript; its released identity
and revision note are retained.
\item[Reason] The continuum-recorder and ontology refinements need explicit
consumer references in the coupling/completion/stabilization account.
\item[Resolution] A concise source interface states the compiler hypotheses
and operational non-entailment result with their contextual owners.
\item[Retained] The local kernel and contraction theorems, ordinary physical
measurement, and exact canonical binary recorder remain unchanged.
\item[Open boundary] Physical continuum source selection, universal apparatus
control, and objective actualization remain separate research obligations.
\end{description}

\section*{Version 6 Revision Note}
\addcontentsline{toc}{section}{Version 6 Revision Note}
\begin{description}[style=nextline]
\item[Supersedes] \emph{Measurement as Disturbance and Stabilization in Modal
Triplet Theory}, version 5.
\item[Reason] The earlier paper combined localized disturbance, global
contractivity, multiple outcome basins, decoherence, Born probabilities,
uncertainty, Bell correlations, and neutrino mass in one mechanism. A global
contraction has only one fixed point, basin volume does not by itself produce
Born weights, and noninjective projection does not imply that no right inverse
or representative section can exist.
\item[Resolution] This version separates coupling, transition completion,
outcome-conditioned stabilization, decoherence, probability source, and
one-history actualization. It adds the missing completion kernel, corrects the
inverse terminology, and states the current q79 recorder result at its exact
restricted tier.
\item[Retained result] A measuring apparatus can be treated as a localized
physical coupling followed by basin-local stabilization. Contractive dynamics
can explain record persistence and repeatability after an outcome has been
resolved.
\item[Remaining boundary] MTT has not yet derived the completion instrument
and Born-compatible source law for every apparatus context, nor an objective
rule selecting one ontic history. The physical upper geometry must be supplied
by the selected branch rather than assumed to be a generic ten-dimensional
product.
\end{description}

\tableofcontents

\section{The measurement chain}

A measurement is not important because a conscious observer notices it. It is
important because a physical interaction produces a durable, communicable
record. A complete model must answer several different questions:
\begin{enumerate}
\item How does the apparatus couple to the system?
\item Which record alternatives are physically available?
\item What assigns a probability to each alternative?
\item What state is available after each recorded outcome?
\item Why does a completed record remain stable?
\item If an ontic account is intended, what selects one individual history?
\end{enumerate}

The corrected MTT decomposition is
\[
\boxed{
\text{coupling/disturbance}
\longrightarrow
\text{outcome completion}
\longrightarrow
\text{record stabilization}.
}
\]
The middle arrow cannot be omitted. A disturbance may move a state toward a
boundary, and a contraction may stabilize a state once it lies in a basin.
Neither statement determines which basin receives the state.

\section{Standard operational target}

Let \(\cH\) be a complex Hilbert space and \(\rho\) a density operator.

\begin{definition}[Quantum instrument]
A finite quantum instrument
\(\cI=\{\cI_i\}_{i\in I}\) is a family of completely positive,
trace-nonincreasing maps on trace-class operators such that
\(\sum_i\cI_i\) is trace preserving. Its effects are
\[
E_i=\cI_i^*(\id),\qquad
E_i\geq0,\qquad
\sum_iE_i=\id.
\]
\end{definition}

The outcome probability and conditional state are
\begin{equation}
p_i=\Tr[\cI_i(\rho)]=\Tr(\rho E_i),
\qquad
\rho_i'=\frac{\cI_i(\rho)}{p_i}
\quad(p_i>0).
\label{eq:instrument}
\end{equation}
An instrument therefore includes both the classical record \(i\) and the
state passed to later physical interactions \cite{DaviesLewis1970}.

The nonselective channel
\[
\cI_{\mathrm{ns}}=\sum_i\cI_i
\]
describes what remains if the outcome label is ignored. It is not an
outcome-selection law. Keeping \(\cI_{\mathrm{ns}}\) while deleting the
individual \(\cI_i\) removes precisely the information needed to say which
record occurred.

\section{An upper transition-completion kernel}

MTT seeks an upper physical account of equation \eqref{eq:instrument}. The
minimal classical-measure analogue is an outcome-resolved kernel.

\begin{definition}[Transition-completion kernel]
Let \(E\) be a measurable pre-completion exit space, let \(I\) be a finite
record set, and let \(D_i\) be the post-completion state space associated with
record \(i\). A transition-completion kernel is a family
\[
\cK_i:E\times\mathfrak{B}(D_i)\longrightarrow[0,1]
\]
such that:
\begin{enumerate}
\item for fixed \(x\in E\), \(\cK_i(x,\cdot)\) is a finite measure on \(D_i\);
\item for fixed measurable \(A\subseteq D_i\),
\(\cK_i(\cdot,A)\) is measurable; and
\item for every \(x\in E\),
\[
\sum_{i\in I}\cK_i(x,D_i)=1.
\]
\end{enumerate}
\end{definition}

The kernel can encode unresolved upper degrees of freedom, an effective
stochastic closure, or a genuinely probabilistic law. Those interpretations
are not equivalent. The definition only states the data required at the
effective level.

\begin{theorem}[Completion law]
\label{thm:completion}
Let \(\nu\) be a probability measure on \(E\) and
\(\{\cK_i\}_{i\in I}\) a transition-completion kernel. Define
\begin{equation}
p_i=\int_E\cK_i(x,D_i)\,d\nu(x).
\label{eq:completion-prob}
\end{equation}
Then \(p_i\geq0\) and \(\sum_i p_i=1\). If \(p_i>0\), the conditional
post-completion law
\begin{equation}
\nu_i'(A)
=\frac{1}{p_i}\int_E\cK_i(x,A)\,d\nu(x),
\qquad A\in\mathfrak{B}(D_i),
\label{eq:completion-state}
\end{equation}
is a probability measure on \(D_i\).
\end{theorem}

\begin{proof}
Nonnegativity follows from positivity of each kernel measure. Finite
additivity and normalization give
\[
\sum_i p_i
=\int_E\sum_i\cK_i(x,D_i)\,d\nu(x)
=\int_E1\,d\nu(x)=1.
\]
For \(p_i>0\), equation \eqref{eq:completion-state} inherits countable
additivity from \(\cK_i(x,\cdot)\), is nonnegative, and satisfies
\(\nu_i'(D_i)=1\).
\end{proof}

\begin{remark}
Theorem~\ref{thm:completion} is a normalization theorem, not a source theorem.
It says what follows once \(\nu\) and \(\cK\) have been physically selected.
It does not determine either object and does not imply Born weights.
\end{remark}

A deterministic completion is the special case in which a map
\(c:E\to\bigsqcup_iD_i\) sends every exit state to one outcome-conditioned
state:
\[
\cK_i(x,A)
=\mathbf{1}_{\{c(x)\in D_i\cap A\}}.
\]
An ensemble of hidden upper states can then yield nontrivial effective
probabilities through \(\nu\), even though \(c\) is deterministic. This is an
ordinary pushforward mechanism. It must not be described as probability
created by projection.

\section{Basin-local stabilization}

Once completion has produced a state in \(D_i\), an outcome-conditioned map
may stabilize the corresponding record.

\begin{theorem}[Record-basin contraction]
\label{thm:basin-contraction}
Let \((D_i,d_i)\) be nonempty and complete, and let
\(T_i:D_i\to D_i\) satisfy
\[
d_i(T_i x,T_i y)\leq q_i d_i(x,y),
\qquad 0\leq q_i<1.
\]
Then \(T_i\) has a unique fixed record \(r_i\in D_i\), and
\[
d_i(T_i^n x,r_i)\leq q_i^n d_i(x,r_i)
\]
for every \(x\in D_i\).
\end{theorem}

\begin{proof}
This is the Banach contraction theorem applied on the explicitly declared
complete invariant basin \(D_i\).
\end{proof}

The hypotheses matter. Completeness, nonemptiness, invariance
\(T_i(D_i)\subseteq D_i\), and the strict contraction constant are all proof
obligations. A local linear estimate near a candidate record does not
establish a global invariant basin.

\begin{proposition}[A global contraction cannot encode several outcomes]
\label{prop:no-global-contraction}
If a map \(T:D\to D\) is a contraction on one nonempty complete space \(D\),
then it cannot have two distinct fixed records.
\end{proposition}

\begin{proof}
The Banach theorem gives a unique fixed point. Equivalently, if
\(Tx=x\) and \(Ty=y\), then
\[
d(x,y)=d(Tx,Ty)\leq qd(x,y)
\]
with \(q<1\), hence \(d(x,y)=0\).
\end{proof}

This corrects a central tension in version 5. Multiple outcomes require
separate invariant basins \(D_i\), a noncontractive transition region, a
context-dependent map, an outcome kernel, or some combination of these. A
single global Fundamental Contractivity Condition cannot simultaneously
select several stable records.

\subsection{Repeatability}

Repeatability is a statement about what happens after record \(i\) is
completed. If the later apparatus interaction preserves \(D_i\) and its
readout identifies the same record throughout a neighborhood of \(r_i\), then
Theorem~\ref{thm:basin-contraction} supplies asymptotic stabilization.
Immediate exact repeatability requires a stronger nondemolition or idempotence
condition on the instrument. It does not follow from attraction alone.

\section{Projection, decoding, and irreversibility}

Let \(P:\cU\to\cX\) be an effective projection from upper configurations to
records. If \(P\) is noninjective, two upper states \(u_1\neq u_2\) may have
the same record \(P(u_1)=P(u_2)\).

\begin{proposition}[Unique decoding versus representative selection]
\label{prop:decoder}
If \(P:\cU\to\cX\) is noninjective, there is no decoder
\(D:\cX\to\cU\) satisfying
\[
D\circ P=\operatorname{id}_{\cU}.
\]
A right inverse or section \(s:\cX\to\cU\) satisfying
\[
P\circ s=\operatorname{id}_{\cX}
\]
may nevertheless exist. Such a section chooses one representative and does
not recover the actual upper state in a nontrivial fiber.
\end{proposition}

\begin{proof}
If \(P(u_1)=P(u_2)\), a decoder would imply
\[
u_1=D(P(u_1))=D(P(u_2))=u_2,
\]
contradicting noninjectivity. The second statement has no such contradiction:
a section chooses one element of each represented fiber.
\end{proof}

Irreversibility must therefore be stated operationally. Possible precise
claims include:
\begin{itemize}
\item failure of unique decoding of the actual upper history;
\item merger of distinct effective histories under the record map;
\item absence of a physically admissible recovery channel;
\item hysteresis or entropy production in a declared open-system model;
\item or failure of upper dynamics to descend to a reversible lower map.
\end{itemize}
Bare noninjectivity does not prove all of them.

\section{Decoherence and outcome completion}

Suppose a pointer decomposition defines a dephasing channel
\[
\cD(\rho)=\sum_iP_i\rho P_i.
\]
Decoherence controls the off-diagonal blocks of the reduced state and explains
why interference between record alternatives can become negligible
\cite{Zurek2003}. It does not, by itself, choose an index \(i\).

The distinction is visible algebraically:
\[
\rho
\longmapsto
\cD(\rho)
=\sum_iP_i\rho P_i
\]
is a nonselective channel. An outcome-resolved instrument retains the
individual maps
\[
\cI_i(\rho)=P_i\rho P_i
\]
and associates each with a record. Even then, the instrument gives an ensemble
law and conditional states. If the theory promises one objective ontic
history, another source or selection statement is required.

In MTT language, decoherence may occur inside or between approximate record
sectors, while basin completion says which record label is emitted. Record
stabilization then keeps that label robust. These are adjacent physical
processes, not one theorem.

\section{The Born source obligation}

Given an instrument, standard quantum mechanics predicts
\[
p_i^{\mathrm{QM}}=\Tr(\rho E_i).
\]
Given an upper completion package, Theorem~\ref{thm:completion} predicts
\[
p_i^{\mathrm{upper}}
=\int_E\cK_i(x,D_i)\,d\nu_\rho(x).
\]
The required MTT equality is
\begin{equation}
\boxed{
\int_E\cK_i(x,D_i)\,d\nu_\rho(x)
=\Tr(\rho E_i)
\quad\text{for every allowed }\rho,\cI,i.
}
\label{eq:born-source}
\end{equation}

Equation \eqref{eq:born-source} displays all the missing data:
\begin{itemize}
\item the preparation-dependent upper law \(\nu_\rho\);
\item the selected apparatus completion kernel \(\cK\);
\item the quantum effect \(E_i\);
\item and the equality on a declared preparation and apparatus domain.
\end{itemize}

Relative basin volume
\[
\frac{\mu(B_i)}{\sum_j\mu(B_j)}
\]
is merely a normalized probability model until the measure, preparation
dependence, and equality \eqref{eq:born-source} are proved. Contractivity does
not determine those weights. The version 5 ``Born rule from basin measures''
proved only that a random initial point lands in a basin with the measure
assigned to that basin; it did not prove equality to quantum trace weights.

\section{Current q79 measurement status}

\subsection{Exact canonical domain}

The q79 program now supplies a selected operational realization on one
restricted domain. For the canonical binary one-anchor nondemolition Fock
recorder:
\begin{itemize}
\item the finite state, observable, and output-algebra data are fixed;
\item the commuting output algebra defines the record events;
\item the selected normal state emits the stopped output measure;
\item second-moment capture descent is exact; and
\item no separate Born axiom, stochastic primitive, observed probability, or
fit is inserted on that domain.
\end{itemize}

Thus the analogue of equation \eqref{eq:born-source} is closed for that
canonical binary apparatus. This is stronger than the conditional basin-volume
story in version 5.

\subsection{Open quantifiers}

\label{sec:recorder-source-interface}
The Born/record companion explains the selected Fock law and its conditional
continuum compiler \cite{NeroBornClassical2026,FrozenFock,FrozenContinuum}.
The compiler requires a supplied positive self-adjoint Hessian, an invariant
rank-three sector with kernel/support ranks one and two, and an isometry
intertwining its projector pair with the finite pair. With the same clock,
minimal Luders coupling, and no added Hamiltonian, it transports both the
stopped probabilities and conditional states. Approximate intertwining gives
finite-horizon error bounds; normalized rare-event states require a positive
weight floor. This is how a selected geometric source could supply this
paper's completion instrument. It does not select the physical continuum
endpoint, finite comparison map and tails, or clock by itself.

The ontology source has a different consumer role
\cite{LocalityCompanion,FrozenOntology}. Its 448-atom canonical checkpoint
admits one record per atom with the same probabilities and conditional states
as the operational instrument; a coactual completion can preserve those data
as well. Thus the completed operational instrument does not entail either
ontology. This context-specific countermodel is not an apparatus-independent
hidden-variable model, a physical one-history selector, or an additional
stabilization law. In particular, stabilization after a record cannot choose
between completions that already agree on every retained operational datum.

The exact result does not yet establish:
\begin{itemize}
\item the completion instrument for every allowed preparation and apparatus;
\item finite-bandwidth and non-Markov detector corrections;
\item a contextual family of overlapping measurements;
\item a universal basin geometry for all quantum observables;
\item or objective selection of one ontic history.
\end{itemize}

The correct ledger is
\[
\begin{array}{ll}
\text{canonical q79 binary recorder:}&\text{exact on its domain},\\
\text{general MTT apparatus family:}&\text{open},\\
\text{universal Born source theorem:}&\text{open},\\
\text{objective one-history selector:}&\text{open}.
\end{array}
\]

\section{A guarded Ornstein--Uhlenbeck model}

Version 5 used Ornstein--Uhlenbeck (OU) widths as an explanation of quantum
uncertainty. The OU calculation is valid as a conditional linear-response
model, but the identification with Heisenberg uncertainty is not derived.

Let \(W_t\) be standard Brownian motion and let a square-integrable \(a_0\)
be independent of its future increments. Suppose a reduced disturbance
coordinate satisfies
\begin{equation}
da_t=-\gamma a_t\,dt+\sigma\,dW_t,
\qquad \gamma>0.
\label{eq:ou}
\end{equation}

\begin{lemma}[OU stationary variance]
\label{lem:ou}
Equation \eqref{eq:ou} has the solution
\[
a_t=e^{-\gamma t}a_0
+\sigma\int_0^t e^{-\gamma(t-s)}\,dW_s,
\]
and
\[
\Var(a_t)
=e^{-2\gamma t}\Var(a_0)
+\frac{\sigma^2}{2\gamma}
\left(1-e^{-2\gamma t}\right).
\]
Consequently its stationary variance is
\[
\Var_{\mathrm{stat}}(a)=\frac{\sigma^2}{2\gamma}.
\]
\end{lemma}

\begin{proof}
The variation-of-constants formula gives the displayed solution. Ito isometry
gives
\[
\sigma^2\int_0^t e^{-2\gamma(t-s)}\,ds
=\frac{\sigma^2}{2\gamma}(1-e^{-2\gamma t}),
\]
which proves the variance formula and its limit.
\end{proof}

If instead
\[
\dot a=-\gamma a+\eta(t),
\qquad
\|\eta\|_\infty\leq M,
\]
then variation of constants gives
\[
\limsup_{t\to\infty}|a(t)|\leq\frac{M}{\gamma}.
\]
The stochastic variance and deterministic amplitude bound have different
meanings and should not be identified.

To turn Lemma~\ref{lem:ou} into a quantum-uncertainty theorem, MTT would need
to derive:
\begin{enumerate}
\item the reduced coordinate \(a\) from selected observables;
\item the noise or unresolved-state law and its strength \(\sigma\);
\item the damping \(\gamma\) from the same source;
\item the canonical commutator or symplectic form; and
\item the Robertson--Schrodinger lower bound with the correct \(\hbar\)
normalization.
\end{enumerate}
Without those steps, the OU floor is an effective fluctuation model, not the
origin of quantum uncertainty.

\section{A concrete apparatus example}

Consider a Stern--Gerlach-type spin readout. The corrected stages are:
\begin{enumerate}
\item A magnetic-field gradient couples spin and position, producing
spatially distinguishable wave packets.
\item Environmental coupling and detector amplification suppress
interference between macroscopic record sectors.
\item An outcome-resolved instrument assigns the two detector records and
their conditional post-measurement states.
\item Record-sector dynamics stabilizes the fired detector state.
\end{enumerate}

The field gradient is the physical coupling. Decoherence helps make the
records robust. The instrument supplies the outcome-resolved law.
Stabilization explains persistence. Calling the first step a disturbance does
not calculate the third step, and calling the fourth step a contraction does
not select between the two detector channels.

The same decomposition applies to interferometry, weak measurements, and
multi-stage protocols. Bell experiments additionally require a bipartite state,
local instrument algebras, setting assumptions, and a locality analysis. They
cannot be derived from stabilization alone.

\section{Relation to interpretations}

The framework is compatible with several ontological readings unless stronger
MTT data are added:
\begin{itemize}
\item As an operational theory, the instrument and output law suffice for
record statistics.
\item As an epistemic upper-state model, probabilities may reflect a
preparation law over unresolved upper configurations.
\item As a stochastic completion, the kernel may be primitive or derived from
a limit theorem.
\item As a deterministic ontic completion, a selected upper state and
deterministic completion map must explain the observed statistics.
\end{itemize}

No observer-dependent collapse is required to formulate any of these options.
But the options are physically different. A claim that MTT selects one of
them must identify the relevant source theorem.

\section{Completion program}

The next measurement theorem should be built in this order:
\begin{enumerate}
\item Select an upper carrier, preparation map, and apparatus coupling from
the physical q79 branch.
\item Identify finite record algebras and outcome spaces.
\item Derive every instrument element or transition-completion kernel from
the same coupling.
\item Prove equation \eqref{eq:born-source} for a nontrivial family of
preparations and apparatus contexts.
\item Prove basin-local stabilization or nondemolition repeatability after
completion.
\item Derive controlled detector bandwidth, memory, and environmental
corrections.
\item State whether one-history actualization is outside the theory or supply
its selected mechanism.
\end{enumerate}

This order prevents three familiar substitutions:
\[
\begin{gathered}
\text{decoherence}\neq\text{outcome selection},\\
\text{basin stability}\neq\text{Born weight},\\
\text{projection}\neq\text{probability source}.
\end{gathered}
\]

\section{Conclusion}

The disturbance-and-stabilization intuition survives, but only as two parts of
a three-stage measurement process. A physical coupling can displace a state.
An outcome-resolved kernel or instrument completes that interaction into a
record and post-record state. Basin-local contraction can then stabilize the
record. A global contraction cannot perform all three jobs because it has only
one fixed point.

The corrected inverse language also matters. Noninjective projection forbids
unique decoding of the actual upper state, but it does not forbid every right
inverse or representative section. Irreversibility requires a specified
physical decoder, channel, or dynamical criterion.

MTT has one substantial exact foothold: the canonical q79 binary one-anchor
recorder emits its stopped-output law and exact capture descent on the declared
domain without a fitted probability. The general apparatus and one-history
problems remain open. The transition-completion kernel introduced here makes
their missing content explicit and gives the measurement program a testable,
noncircular next step.

% BEGIN MTT MANAGED COMPUTATIONAL EVIDENCE
\section*{Computational Evidence and Reproducibility}
The distinction among coupling, outcome completion, and record stabilization is established from instruments and transition kernels. The open strict-upgrade ledger does not select an outcome law and is cited only to mark the stronger unresolved source boundary.

The referenced rows are frozen to the curated results repository state identified below. Hashes are grouped in eight-character blocks for line breaking.
\begin{quote}\small
\textbf{Repository:} \url{https://github.com/PeterNero/mtt-results-repro}\\
\textbf{Commit:} \texttt{31247ebb 5c22f3fb b5443024 365433c6 ee0bff4a}\\
\textbf{Manifest:} \path{release/result_manifest.json}\\
\textbf{Manifest SHA-256:}\\
\texttt{fb399689 60b00584 631dbf53 1a708e18}\\
\texttt{ef928d6b 6d935119 c185d7f6 32b1e7cd}
\end{quote}

Tier labels are quoted verbatim from that manifest. A row used directly supports only the specific computational statement identified above; a corpus-state cross-check does not prove this paper's local theorems; and an open row is evidence of an unresolved obligation, never of closure.

\par\begingroup\dimen0=\pagegoal\advance\dimen0 by -\pagetotal\ifdim\dimen0<7\baselineskip\newpage\fi\endgroup
\paragraph{Open boundary (not evidence of closure).}
\begin{itemize}
\setlength{\itemsep}{0.25em}
\item \path{A05/strict_upgrade_ledger} (\emph{open}).\par Historical 2/9
strict no-knob ledger snapshot, not a current global completion count.
\end{itemize}

No imported row changes theorem ownership or promotes a neighboring claim: all local statements retain their stated hypotheses, domains, and limitations.
% END MTT MANAGED COMPUTATIONAL EVIDENCE

\begin{thebibliography}{99}
\raggedright

\bibitem{DaviesLewis1970}
E.~B. Davies and J.~T. Lewis,
\emph{An Operational Approach to Quantum Probability},
Communications in Mathematical Physics \textbf{17} (1970) 239--260,
doi:10.1007/BF01647093.

\bibitem{Zurek2003}
W.~H. Zurek,
\emph{Decoherence, Einselection, and the Quantum Origins of the Classical},
Reviews of Modern Physics \textbf{75} (2003) 715--775,
doi:10.1103/RevModPhys.75.715,
arXiv:quant-ph/0105127.

\bibitem{NeroBornClassical2026}
P.~Nero,
\emph{Born-Compatible Record Measures and the Classical Concentration Limit:
Separate Theorems and Their MTT Interface},
version 3, MTT research manuscript, September 2026, Section 4.

\bibitem{NeroContextOrder2026}
P.~Nero,
\emph{Contextuality and Sequential Measurement Order: Distinct Obstructions
with a Shared MTT Interface},
version 2, MTT research manuscript, July 2026.

\bibitem{MTTResults}
P.~Nero,
\emph{MTT Results Reproducibility Repository},
\url{https://github.com/PeterNero/mtt-results-repro}.

\bibitem{FrozenFock}
P.~Nero, \emph{Canonical q79 Fock output measure and second-moment capture descent}, frozen source (2026).
\url{https://github.com/PeterNero/mtt-results-repro/blob/f141a20ea23c5c3ff19cc2161c0e226e29ade8a7/release/results/q79_fock_output_measure/artifact.json}.

\bibitem{FrozenContinuum}
P.~Nero, \emph{Continuum Hessian-to-recorder compiler}, frozen source (2026).
\url{https://github.com/PeterNero/mtt-results-repro/blob/f141a20ea23c5c3ff19cc2161c0e226e29ade8a7/release/results/q79_continuum_recorder_compiler/artifact.json}.

\bibitem{FrozenOntology}
P.~Nero, \emph{q79 operational ontology non-entailment}, frozen source (2026).
\url{https://github.com/PeterNero/mtt-results-repro/blob/f141a20ea23c5c3ff19cc2161c0e226e29ade8a7/release/results/q79_ontology_nonentailment/artifact.json}.
\bibitem{LocalityCompanion}
P.~Nero, \emph{Locality, Coherent Alternatives, and Physical Records:
An Interpretive Account of Quantum Experiments in Modal Triplet Theory},
unpublished version 2 (September 2026), subsection \emph{Operational data do not
force many actual worlds}.

\end{thebibliography}

\end{document}
